% Generated draft: requires mathematical and source review.
\documentclass[11pt]{article}
\usepackage[margin=25mm]{geometry}
\usepackage{fontspec}
\setmainfont{Latin Modern Roman}
\usepackage{amsmath,amssymb,mathtools,unicode-math}
\setmathfont{Latin Modern Math}
\usepackage{enumitem,needspace,xurl,hyperref}
\setlength{\parindent}{0pt}
\setlength{\parskip}{5pt}
\setlength{\emergencystretch}{3em}
\newcommand{\N}{\mathbb{N}}
\newcommand{\Z}{\mathbb{Z}}
\newcommand{\Q}{\mathbb{Q}}
\newcommand{\R}{\mathbb{R}}
\newcommand{\C}{\mathbb{C}}
\newcommand{\F}{\mathbb{F}}
\newcommand{\A}{\mathbb{A}}
\newcommand{\PP}{\mathbb{P}}
\newcommand{\E}{\mathbb{E}}
\newcommand{\eps}{\varepsilon}
\newcommand{\dd}{\,\mathrm{d}}
\begin{document}
\section*{UnsolvedMath ID: 4200011}
\textbf{Source problem code:} catalogue entry
\subsection{Definitions and mathematical statement}
In the surface setting, let $M$ be a connected closed smooth surface with a smooth area probability measure $\mu$. Write $\operatorname{Diff}^1_\mu(M)$ for its area-preserving $C^1$ diffeomorphisms with the $C^1$ topology. A property is generic in the Baire-category sense if it holds on a residual subset, namely a countable intersection of open dense subsets.

A diffeomorphism is uniformly hyperbolic (Anosov) if there is a continuous invariant splitting $TM=E^s\oplus E^u$ and constants $C>0$, $0<\lambda<1$ such that, for all $n\geq0$,
\[
\|Df^n v\|\leq C\lambda^n\|v\|\quad(v\in E^s),\qquad
\|Df^{-n}v\|\leq C\lambda^n\|v\|\quad(v\in E^u).
\]
The standard surface dichotomy identified by the supplied status evidence is that there exists a residual $\mathcal R\subseteq\operatorname{Diff}^1_\mu(M)$ such that every $f\in\mathcal R$ is Anosov or has all Lyapunov exponents zero at $\mu$-almost every point. Lyapunov exponents refer to the asymptotic logarithmic growth rates of the singular values of $Df^n_x$, where they exist. Zero almost everywhere does not mean zero at every point.

The historical source asks for a generic hyperbolic-or-zero dichotomy on a ``compact manifold''. Beyond dimension two, area-preserving must be replaced or interpreted, for example as preserving a smooth volume or a symplectic form, and hyperbolic must be distinguished from nonuniform hyperbolicity. The source does not fix these choices. Its supplied assessment says the surface case is solved and that higher-dimensional dynamics has weaker dominated/partially-hyperbolic alternatives; it simultaneously describes the naive higher-dimensional dichotomy as failing or remaining open. Thus no single precise unresolved higher-dimensional theorem is determined by this record.
\subsection{Short English statement}
For $C^1$-generic area-preserving surface diffeomorphisms, the supplied status evidence identifies the Anosov-or-zero-Lyapunov-exponents-almost-everywhere dichotomy as solved. The source's higher-dimensional extension needs a precise measure class and a revised alternative.
\subsection{Sources}
\begin{itemize}
\item Local UnsolvedMath registry, permanent ID 4200011, exact source code catalogue entry, ``Ma\~n\'e's last theorem''; attributed to Ricardo Ma\~n\'e, proposed year 1995. Registry status partially solved, classification PARTIAL-PROGRESS, with higher-dimensional cases singled out.
\item Oliver Knill, \emph{Open problems in Hamiltonian dynamics} (2000), item 11, ``MA\~N\'E'S LAST THEOREM'', \url{https://people.math.harvard.edu/~knill/seminars/intr/}. Supplied HTML extraction accessed July 29, 2026.
\item Bochi, \texttt{arXiv:math/0202233} (2002), cited in the supplied status evidence as proving the hyperbolic-or-zero dichotomy for compact surfaces. No full publication citation is supplied. The higher-dimensional dominated/partially-hyperbolic progress is reported without precise references or a specific replacement theorem.
\item Supplied provenance flags rights NEEDS\_REVIEW: public access was verified, but redistribution permission was not established. The surface result and higher-dimensional qualifications are retained as supplied status evidence, not an independent literature check.
\end{itemize}
\end{document}
